\documentclass[a4paper]{article}

\usepackage{graphicx}
\usepackage{amsmath,amssymb}
\usepackage{minted}
\usepackage{multirow}
\usepackage{float}
\usepackage{todonotes}
\usepackage{forest}
\usepackage{xstring}
\usepackage{xcolor}
\usepackage[most]{tcolorbox}
\usepackage{xparse}

\usetikzlibrary{graphs,graphdrawing}
\usegdlibrary{layered}

% for external graphviz
\input{/home/donkarlo/Dropbox/repo/nd_language_ai_project/src/nd_language_ai/formal/computer/programming/tex/latex/code_clips/external_graphviz.tex}

%for tags, comma separated \semtag{tag1, tag2}
\input{/home/donkarlo/Dropbox/repo/nd_language_ai_project/src/nd_language_ai/formal/computer/programming/tex/latex/parts/control_sequence/kind/semtag/semtag.tex}

% for links, just type \seehere
\input{/home/donkarlo/Dropbox/repo/nd_language_ai_project/src/nd_language_ai/formal/computer/programming/tex/latex/code_clips/seehere.tex}

\title{}
\date{}

\begin{document}
    \maketitle
    
    
    \section{Nun}
        → wird benutzt um einen Gedanken weiterzuführen, zu reagieren oder eine Schlussfolgerung zu machen (ähnlich wie well, so, then).
        Beispiel:
        Nun, das ist ein Problem.
        Well, that is a problem.
        
        Auch es Verstärkend in Fragen\\
        → oft mit was, wie, wer usw.
        
        Beispiel:\\
        Was sollen wir nun tun?\\
        What should we do now?
        
        Verschiden swichen jezt un nun\\
        Wichtiger Unterschied\\
        jetzt → neutrales „now“\\
        nun → etwas literarischer oder formeller und oft logische Weiterführung im Text
    % \bibliographystyle{plain}
    % \bibliography{/home/donkarlo/Dropbox/repo/nd_shared_research/bibliography/bibliography.bib}
\end{document}